% Source: 26_Chapter_17_God_s_Consul_Gregory_the_Great.txt
\bookchapter[God’s Consul]{CHAPTER 17}{GOD’S CONSUL}{Gregory the Great}

\leadin{E}{ARLY IN AD 590} Rome was in agony. The city suffered through the tragedies of floods and
the atrocities of war only to be smitten by the relentless spread of the plague. Men felt hardly
more than a little soreness of the throat; afterward came the black eruptions and a swift death. The
carts were piled high with corpses. People went insane. Rome became a desert, and the pope himself,
Pelagius II, died, screaming in agony.

For six months no pope ruled in St. Peter’s Basilica. When church leaders decided to elect a monk
named Gregory, he refused the office and even fled from the city, hiding in the forest, until he was
found and dragged back to Rome. After notifying Constantinople, officials consecrated him St.
Peter’s successor on September 3, 590.

Gregory was a most unlikely candidate for greatness. Fifty, balding, and frail, he had no craving
for the papal office. He complained that he was “so stricken with sorrow that he could scarcely
speak.” Yet he began his administration with a public act of humiliation because the plague had
taken the life of his predecessor. Seven processions filed through the streets for three days.
Prayers were said; hymns were sung. But to no immediate avail. The plague continued to ravage the
city. Then, mercifully, it seemed to subside.

A later legend traced the staying of the calamity to Gregory’s action and told of the appearance of
the archangel Michael, who put back his drawn sword into its sheath over the mausoleum of the
emperor Hadrian. Since that time Romans have called it the Castle of St. Angelo. They adorned it
with the statue of an angel. Tourists can still find it on the banks of the Tiber.

Rome was a symbol of the continent. What we now call Europe arose like a phoenix from the blazing
ruins of a devastated empire. And more than any other force, it was Christianity that brought life
and order out of the chaos. How? What did Christianity bring to the devastation to erect a new order
called Christian Europe?

\section{GREGORY’S PLACE IN HISTORY}

The church enlisted the Celtic and Benedictine monks to serve as a spiritual militia for winning the
barbarians to the Christian faith. It turned to the papal office to provide some stable structure
for a new way of life. And it read and recounted the ideals of Augustine’s theology to maintain a
framework of spiritual meaning. No man mastered these instruments of the future better than Gregory.

In his book Pastoral Care, Gregory stressed that the spiritual leader should never be so absorbed in
external cares as to forget the inner life of the soul, nor neglect external things in the care for
his inner life. “Our Lord continued in prayer on the mountain,” Gregory wrote, “but wrought miracles
in the cities; showing to pastors that while aspiring to the highest, they should mingle in sympathy
with the necessities of the infirm. The more kindly charity descends to the lowest, the more
vigorously it recurs to the highest.”

The words were autobiographical.

When Gregory died in 604, worn out after thirty years of prayers on the mountains and miracles in
the cities, his epitaph proclaimed him “God’s Consul.” It was a singularly appropriate description
of the man who had exerted himself to the utmost to be solely God’s while ruling church and world
like a Roman statesman, the last of his line.

Not long after his death the churches came to speak of him as “Gregory the Great.” And in time the
Catholic Church added his name to those of Augustine, Ambrose, and Jerome to speak of the “Latin
Fathers of the Church.” In terms of intellectual powers alone, Gregory probably doesn’t belong in
such company. But he combined great executive ability with warm sympathy for human need. And if
goodness is the highest kind of greatness, then the church moved rightly in according him the title
“Great.” Certainly no other man or woman better represents the early Middle Ages. Gregory was born
about AD 540 from an old and wealthy senatorial family of Rome and was educated for government
service. He could hardly have stepped on the pages of history at a more angry time. In his childhood
Rome changed hands over and over again. He was fourteen in the year 554 when Narses became viceroy
of Italy under the Emperor Justinian in Constantinople. Then at last, with the Visigoth rule over
Italy destroyed, a few brief years of peace followed before the savage Lombards began their campaign
of burning churches, slaying bishops, robbing monasteries, and reducing cultivated fields to a
wilderness. Rome was clearly no longer the metropolis Ambrose and Augustine had known. The city of
the caesars was fast becoming the city of the popes; and it was Gregory’s fate to appear at the hour
of transition.

Suddenly, at the age of 33, Gregory found that the Emperor Justin had appointed him prefect (or
mayor) of Rome, the highest civil position in the city and its surrounding territory. The whole
economy of Rome—the grain supplies, the welfare program for the poor, the construction of buildings,
baths, sewers, and riverbanks—rested on Gregory’s shoulders. To make the burden even heavier, his
appointment in 573 came just as both the pope and Narses died.

Gregory was never comfortable with worldly power, however. He preferred the solitude of a monastic
cell. Within a few years he stepped down from public office and broke with the world. Upon the death
of his father, he spent the greater part of his personal fortune in founding seven monasteries. He
distributed the rest in alms for the poor, then laid aside all vestiges of rank and transformed his
father’s palace into a monastery dedicated to St. Andrew. He exchanged the purple toga for the
coarse robe of a monk and began to live with extraordinary asceticism, eating only raw fruit and
vegetables, praying most of the night, wearing a hair shirt, throwing himself into the many duties
of a Benedictine. He had never been strong, and now unceasing fasting ruined his digestion and
played havoc with his heart. Yet Gregory looked upon these years as the happiest of his life.

Gregory’s gifts, however, could not remain hidden. In 579 Pope Pelagius II made him one of the seven
deacons of the Roman church and sent him as ambassador to the imperial court in Constantinople. His
political training and executive ability fitted him eminently for this post. He returned in 585 and
was appointed abbot of his convent, St. Andrew, but remained available for important public
business.

Gregory was perfectly content to be an abbot, and he would probably have continued stretching for
the other world if the terrible plague had not swept through the city and snuffed out the life of
Pope Pelagius II (579–590).

\section{PICKING UP THE PIECES OF A SHATTERED WORLD}

Soon after Gregory’s election the Lombards laid siege to Rome, and the new pope was forced to cut
short his preaching on Ezekiel: “I am constrained to cease from my exposition, for I am weary of
life. Who can expect me now to devote myself to sacred eloquence, now that my harp is turning to
mourning, and my speech to the voice of them that weep?”

All of Western Europe was in chaos. Serious men, and Gregory was among them, thought that the end of
the world was at hand. “What is it,” he asks in one of his sermons, “that can at this time delight
us in this world? Everywhere we see tribulation, everywhere we hear lamentation. The cities are
destroyed, the castles torn down, the fields laid waste, the land made desolate. Villages are empty,
few inhabitants remain in the cities, and even these poor remnants of humanity are daily cut down.
The scourge of celestial justice does not cease, because no repentance takes place under the
scourge. We see how some are carried into captivity, others mutilated, others slain. What is it,
brethren, that can make us contented with this life? If we love such a world, we love not our joys,
but our wounds.”

The church of Rome survived these attacks, almost the only vestige of organized civilization in the
West that did so. While Gregory regarded his elevation to the papacy as a punishment, he immediately
threw himself into the struggle for order in the midst of chaos.

He wrote urgent letters to the managers of his estates in Sicily: “You sent me a sorry nag and five
good asses: the nag is too wretched to ride, and I simply can’t ride the asses because they are
asses.”

He began a vast correspondence with his bishops: “You evidently paid no attention to my last
letter.”

And he set forth the principles for Christian ministry in his Pastoral Rule: “He who, by the
necessity of his position, is required to speak the highest things, is compelled by the same
necessity to exemplify the highest.”

The prestige of the papacy in the Middle Ages rests in large part on the practical government
maintained by Gregory through these troubled times. He was incessantly busy. Nothing seemed too
great, nothing too little for his personal care. His labors are all the more astonishing when we
consider that he was in poor health and often confined to bed. “For a long time,” he wrote to a
friend in 601, “I have been unable to rise from my bed. I am tormented by the pains of gout; a kind
of fire seems to pervade my whole body: to live is pain; and I look forward to death as the only
remedy.” In another letter he says: “I am daily dying, but never die.”

During Gregory’s time the Church of Rome controlled extensive lands around Rome, in the toe and heel
of Italy and on Sicily. These were called “the patrimony of St. Peter.” Taken together these
estates—something like 1,800 square miles—made the Church the richest landowner in Italy. Quite
naturally, then, when the Lombards invaded central Italy, destroying the imperial administration in
the process, officials of the patrimony stepped in to feed the population and to collect the land
taxes, just as imperial officers had once done. The head of this tax and welfare system was Gregory.

Moreover, when Lombard attacks moved closer and closer to Rome, Gregory undertook the defense of
central Italy. He appointed a military governor and arranged peace with two Lombard leaders. As a
result, after 595, the pope was more important in Lombard politics than any imperial representative.

This participation in the political fortunes of Italy became a significant element for the papal
office in the centuries that followed. After Gregory, the pope was no longer only a Christian
leader; he was also an important political figure in European politics: God’s Consul.

Gregory’s vigorous leadership magnified the authority of the papal chair. One significant clash with
the patriarch of Constantinople reveals his view of the office. In his letters, John IV, the
patriarch, repeatedly used the title universal bishop. It was an honorary title, given to patriarchs
by the emperors Leo and Justinian and confirmed to John and his successors by a Constantinopolitan
synod in 588.

Gregory, however, was provoked and irritated by the assumption of his Eastern rival. He strained
every nerve to get the title revoked. He characterized it as a foolish, proud, profane, wicked,
pestiferous, blasphemous, and diabolical usurpation, and compared anyone who used it to Lucifer. He
threatened to break off communion with the patriarch. He called upon the emperor to punish such
presumption.

In opposition to these high-sounding titles, Gregory called himself simply, “the servant of the
servants of God.” This became one of the standing titles of the popes, although it sounds like irony
when linked with later astounding claims. When a churchman addressed Gregory himself as universal
pope, he strongly repudiated the title, saying, “I have said that neither to me nor to any one else
ought you to write anything of the kind. Away with words which inflate pride and wound charity!”

On the other hand, while Gregory protested high-sounding titles, he claimed and exercised, as far as
he had the opportunity and power, the oversight over the whole church of Christ.

No wonder, therefore, that the successors of Gregory had no scruple in using even more arrogant
titles than the one against which he so solemnly protested with the warning “God resisteth the
proud, but giveth grace to the humble.”

To Gregory pride was a vicious hound that dogged him relentlessly. He speaks of it so often that he
was clearly obsessed with it. He sees pride in all its moods. In his commentary on Job called
Moralia he examines it from every angle: “Pride, which we have called the root of vices, far from
being satisfied with the extinction of one virtue, raises itself up against all the members of the
soul, and as a universal and deadly disease corrupts the whole body.”

It is clear that Gregory himself suffered from the disease. All his frenzy, the continual urgent
letters sent all over Christendom, the determination never to rest for a moment—all these seem to
spring from his knowledge of his own pride, from his desperate desire to face the evil and vanquish
it. Such a lifelong struggle was in perfect harmony with his dedication to monasticism.

\section{MISSIONARY STATESMAN}

With Gregory, monasticism for the first time ascended the papal throne. He continued the austere
simplicity of monastic life, surrounded himself with monks, made them bishops and legates, confirmed
the rule of St. Benedict at a council of Rome, guaranteed the liberty and property of convents, and
by his example and influence rendered the highest service to the monastic life.

A story circulated, probably part fact and part fiction, that Gregory, before he assumed the papal
office, happened upon three young English boys for sale on the Roman marketplace. Gregory was
immediately attracted to them.

“Alas! alas!” he exclaimed, “that such bright faces should be the slaves of an inward darkness. So
beautiful they are and yet their minds are sick and without God’s grace.” He asked the slave owner
what nation they came from. “They are Angles,” answered the man.

“Yes, indeed,” Gregory went on, “they have the appearance of angels and they should be co-heirs of
the angels in Heaven. What province do they come from?”

“From Deira,” said the slave owner, mentioning the ancient name for what is now Northumberland.

“From Deira? Then indeed they should be saved from God’s anger (dei ira) and called to the mercy of
Christ. Who is their King?”

“Aella.”

“Then,” said Gregory, “must Alleluia be sung in Aella’s land.”

Gregory attempted to go to England himself as a missionary monk but was hindered by God and the
pope. But once on the papal throne himself, he sent the Benedictine Augustine and forty monks to
replant the gospel on English soil. The success of that mission in Kent, as we have seen, provides a
direct link for all Anglo-American Christianity with the early church.

\section{DEFENDER OF ORTHODOXY}

Gregory not only yearned to advance the faith in distant places, but he took seriously his calling
as defender of orthodoxy. His teachers in the faith were Ambrose, Augustine, and Jerome, but he
lacked their intellectual abilities. He contributed no new ideas and created no epoch in theology.
But he formulated the common faith of his day and handed it on to the Catholic Church of the Middle
Ages. This faith embraced not merely the official pronouncements of the councils and the teachings
of the Fathers, but also the notions of the illiterate populace, often crude and superstitious, at
times even pagan. To this mass of material he lent the weight of his authority, and as a consequence
it became an integral part of the faith of the Western church, of theologians and bishops as well as
of monks and laymen. It is impossible to follow the thinking of the Middle Ages without the guidance
of Gregory.

In his doctrine of man, Gregory stressed that Adam’s fall affected all his descendants, weakening
but not destroying their freedom of will. Thus once man has been moved by grace, he may cooperate
with it and win merit for himself by his good works, which are the joint product of divine grace and
human will.

In baptism, God grants forgiving grace freely without any merit on man’s part, but for sins
committed after baptism man must make atonement by penance, which is simply a form of punishment
inflicted by the man himself instead of by God: “For either man himself by penance punishes sin in
himself, or God taking vengeance on him smites it.” Penance involves repentance, which must be
sincere and of the heart, and also confession and meritorious works.

The meritorious works, without which penance is not complete, are deeds involving sacrifice or
suffering, such as almsgiving, ascetic practices, and prayers at all hours of the day. The greater
our sins the more we must do to make up for them, and the more careful we must be to avoid them in
the future. Whether we have done enough to atone for them we cannot know until after death.

Fortunately, sinners have the help of the saints. The belief in the intercession of the saints and
the custom of appealing to them to use their influence with Christ did not originate with Gregory;
both the belief and the custom were much older than he. But he emphasized them and made them central
for Christian piety. “Behold,” he wrote, “the severe judge Jesus is about to come; the terror of
that mighty council of angels and archangels is at hand. In that assembly our case will be tried and
yet we are not seeking patrons who will then come to our defense. Our holy martyrs are ready to be
your advocates; they desire to be asked, indeed if I may say so, they entreat that they may be
entreated. Seek them as helpers of your prayer; turn to them that they may protect you in your
guilt.”

Another aid to devotion was the holy relics. Gregory encouraged the collection and veneration of
holy remains of the saints and martyrs: locks of hair, fingernails, toes, pieces of clothing.
Gregory taught—and most of his contemporaries believed—that these items possessed great powers,
including that of self-defense.

And if the saints and relics were not sufficient for righteousness in this life, sins could still be
atoned in purgatory. This is a place of purification and suffering, not for those who die with
serious offenses still charged against them, but for those who are not as yet altogether righteous.
At death the perfectly holy go at once to heaven and the wicked to hell, while those of an
intermediate character, who still have minor sins for which penance must be done, spend a season in
purgatory.

Surely, however, the supreme miracle, the key to all the other expressions of divine power, was the
Holy Eucharist. According to Gregory, the Eucharist is a communion with Christ whose body and blood
are really present in the bread and wine. Feeding upon them we nourish and strengthen our spiritual
life.

The marvelous power of the Eucharist, however, lies in its sacrificial character. It is offered by
the priest for the sins of men—not like Christ’s death upon the cross, for the sins of all men, but
only for the sins of the participants, or of those for whose benefit it may be specifically offered.
For all such it has the same effect as penance, taking the place of a certain amount of suffering
that they would otherwise have to undergo because of their sins. It may benefit the dead as well as
the living, the dead, that is, in purgatory not in hell. If offered for anyone in purgatory, it will
hasten the time of his release.

One of the better examples of Gregory’s belief in the power of the sacred mass is found in his
Dialogues. While the incident is a bit long, it serves as a revealing picture of medieval piety.
After telling of the death of one of his monks who had been found guilty of hoarding money and had
been severely punished, Gregory writes,

\begin{blockquote}
Thirty days later, I began to feel strong compassion for the deceased Justus. As I considered with
deep anguish the penalty he was enduring, I thought of a way to relieve him of his suffering. With
this in mind, I called Pretiosus, the prior, and said to him sadly, “Justus has now been suffering
the torments of fire for a long time and we must show him our charity by helping as much as we can
to gain his release. Beginning today, offer the holy Sacrifice for his soul for thirty consecutive
days. Not one of these days is to pass without a Mass being celebrated for his release.” The prior
obediently accepted the instructions and left.

Days passed, and being busy with other affairs, I lost count of them. Then, one night, Justus
appeared to his brother Copiosus, who asked him at once why he came and how he was. “Up to this
moment I was in misery,” he said, “but now I am well, because this morning I was admitted to
communion.”
\end{blockquote}

Copiosus hurried to tell the monks the good news. Taking exact count of the days, they discovered
that this was the thirtieth consecutive day on which mass had been offered for him. Previous to this
Copiosus did not know that the brethren were offering masses for Justus, nor did the brethren know
that Copiosus had seen him in a vision. At the very moment, therefore, when they became mutually
aware of what had taken place, they realized that the vision and the completion of the thirty masses
occurred at one and the same time. They were now convinced that the brother who had died was freed
from punishment through the Sacrifice of the Mass.

The doctrine was widely accepted in the Western church from Gregory’s time on and helped to give its
peculiar tone to the Christianity of the Middle Ages.

\readinghead{Suggestions for Further Reading}
\begin{readinglist}
\entry Deanesly, Margaret. A History of the Medieval Church 590–1500. London: Methuen, 1969.
\entry Duckett, Eleanor Shipley. The Gateway to the Middle Ages: Monasticism. Ann Arbor: University
  of Michigan Press, 1961.
\entry Gontard, Friedrich. The Chair of Peter. New York: Holt, Rinehart, and Winston, 1964.
\entry *Kardong, Terrence G. Together Unto Life Everlasting: An Introduction to the Rule of
  Benedict. Richardton, ND: Assumption Abbey, 1984.
\entry *Markus, Robert A. Gregory the Great and His World. Cambridge: Cambridge University Press,
  1997.
\entry Zimmerman, Odo John, trans. Saint Gregory the Great: Dialogues. New York: Fathers of the
  Church, 1959.
\end{readinglist}
